\section{Introduction}
\label{sec:intro}

As large language models (LLMs) become central to automated reasoning pipelines, a critical question arises: can these models reliably evaluate the logical validity of their own outputs? Self-evaluation is a cornerstone of LLM-as-judge systems~\citep{zheng2023judging}, self-correction mechanisms~\citep{madaan2023selfrefine}, and safety monitoring frameworks. If models exhibit systematic blindspots when assessing their own reasoning---detecting fallacies less accurately than when evaluating other models' reasoning---the reliability of these systems is fundamentally compromised.

Recent work has established that LLMs struggle with reasoning error detection. \citet{tyen2024llms} showed that GPT-4 achieves only 52.87\% accuracy on mistake finding in chain-of-thought traces, despite succeeding at correction when given the error location. \citet{liu2024self} found self-contradiction rates up to 50\% in certain models, with GPT-4 detecting these contradictions at only 52.2\% F1. Meanwhile, \citet{payandeh2023susceptible} identified a self-detection paradox: GPT-3.5 generates fallacies it cannot subsequently identify. These findings raise a natural question: does \emph{attribution}---whether reasoning is framed as the model's own---systematically bias fallacy detection?

Despite growing interest in LLM self-evaluation, no study has systematically compared detection rates when identical fallacious reasoning is attributed to ``self'' versus ``another model.'' \citet{tyen2024llms} explicitly called for studying cross-model versus self-evaluation differences. \citet{hong2024closer} documented category-specific performance variation across 232 fallacy types but did not examine whether source attribution affects detection. This gap is consequential: if simple attribution framing can bias evaluation, then LLM-as-judge systems and self-correction pipelines may produce systematically skewed assessments.

We address this gap with three controlled experiments using \gptfull and \gptmini on 195 examples from the \logicdataset dataset~\citep{jin2022logical} spanning 13 fallacy categories. Our attribution framing experiment reveals that \gptfull is remarkably robust: the detection gap between self-attributed and other-attributed conditions is only $+3.1\%$ ($p = 0.146$). In genuine cross-model evaluation, self-evaluation accuracy ($82.0\%$) actually exceeds cross-evaluation ($79.5\%$), contradicting the hypothesized self-evaluation penalty. However, we identify category-specific asymmetries reaching $13.3\%$ for false dilemma, faulty generalization, and intentional fallacies.

We make the following contributions:
\begin{itemize}[leftmargin=*,itemsep=0pt,topsep=0pt]
    \item We design and conduct the first controlled study of attribution bias in LLM fallacy detection, isolating the effect of source framing from content differences across 13 fallacy categories.
    \item We perform cross-model evaluation between \gptfull and \gptmini, constructing a $2 \times 2$ detection matrix that reveals self-evaluation slightly outperforms cross-evaluation in genuine settings.
    \item We identify category-specific blindspots---false dilemma, faulty generalization, and intentional fallacies show $+13.3\%$ attribution bias, while fallacy of credibility shows $-13.3\%$ reverse bias---and demonstrate that explicit prompting interventions produce zero improvement ($p = 1.0$).
\end{itemize}
